\begin{thebibliography}{99}\small\raggedright
\bibitem{BM} B. Banaschewski and G. H. Moore, \emph{The dual Cantor--Bernstein theorem and the partition principle}, Notre Dame Journal of Formal Logic \textbf{31} (1990), no.~3, 375--381. \href{https://doi.org/10.1305/ndjfl/1093635502}{doi:10.1305/ndjfl/1093635502}.

\bibitem{Blass} A. Blass, \emph{Injectivity, projectivity, and the axiom of choice}, Transactions of the American Mathematical Society \textbf{255} (1979), 31--59. \href{https://doi.org/10.1090/S0002-9947-1979-0542870-6}{doi:10.1090/S0002-9947-1979-0542870-6}.

\bibitem{BK} A. Blass and D. Kulshreshtha, \emph{Cardinal Well-foundedness and Choice}, arXiv:2310.09643v3, 21 August 2026. \url{https://arxiv.org/abs/2310.09643v3}.

\bibitem{daSilva} S. G. da Silva, \emph{The Axiom of Choice and the Partition Principle from Dialectica Categories}, Logic Journal of the IGPL \textbf{29} (2021), no.~5, 783--797. \href{https://doi.org/10.1093/jigpal/jzaa023}{doi:10.1093/jigpal/jzaa023}.

\bibitem{FO} D. Feldman and M. Orhon, with an appendix by A. Blass, \emph{Generalizing Hartogs' Trichotomy Theorem}, arXiv:0804.0673 (2008). \url{https://arxiv.org/abs/0804.0673}. See Theorem 1 and the discussion of infinite trichotomy.

\bibitem{HT} P. Howard and E. Tachtsis, \emph{No decreasing sequence of cardinals}, Archive for Mathematical Logic \textbf{55} (2016), 415--429. \href{https://doi.org/10.1007/s00153-015-0472-5}{doi:10.1007/s00153-015-0472-5}.

\bibitem{Karagila} A. Karagila, \emph{Embedding orders into the cardinals with $\mathsf{DC}_\kappa$}, Fundamenta Mathematicae \textbf{226} (2014), no.~2, 143--156. \href{https://doi.org/10.4064/fm226-2-4}{doi:10.4064/fm226-2-4}.

\bibitem{OpenAI} OpenAI, \emph{The Partition Principle does not imply Choice}, manuscript dated 24 September 2026, repository revision checked 7 October 2026. Commit \texttt{adc7f1241b42e322a6451854ab7e4b4c146bf78a}. \href{https://github.com/openai/math/tree/adc7f1241b42e322a6451854ab7e4b4c146bf78a/preprints/The-Partition-Principle-does-not-imply-Choice-September-24-2026}{Pinned manuscript and source}.

\bibitem{Pelc} A. Pelc, \emph{On Some Weak Forms of the Axiom of Choice in Set Theory}, Bulletin de l\'Acad\'emie Polonaise des Sciences, S\'erie des Sciences Math\'ematiques, Astronomiques et Physiques \textbf{26} (1978), 585--589. The results used for historical context are discussed in \cite[Section 10]{BK}.

\bibitem{RyanSmith} C. Ryan-Smith, \emph{Local reflections of choice}, Acta Mathematica Hungarica \textbf{176} (2025), no.~1, 244--257. \href{https://doi.org/10.1007/s10474-025-01533-3}{doi:10.1007/s10474-025-01533-3}. Author manuscript: \url{https://academic.calliope.mx/files/papers/2412.13785.pdf}. Section references in this note refer to the checked author manuscript.

\bibitem{Tachtsis2018} E. Tachtsis, \emph{On the comparability of cardinals in the absence of the axiom of choice}, Fundamenta Mathematicae \textbf{242} (2018), no.~3, 247--266. \href{https://doi.org/10.4064/fm376-11-2017}{doi:10.4064/fm376-11-2017}.

\bibitem{Tachtsis2024} E. Tachtsis, \emph{No Decreasing Sequence of Cardinals in the Hierarchy of Choice Principles}, Notre Dame Journal of Formal Logic \textbf{65} (2024), no.~3, 311--331. \href{https://doi.org/10.1215/00294527-2024-0020}{doi:10.1215/00294527-2024-0020}.

\bibitem{FelgnerTruss} U. Felgner and J. K. Truss, \emph{The independence of the Prime Ideal Theorem from the Order-Extension Principle}, J. Symbolic Logic \textbf{64} (1999), 199--215. \href{https://doi.org/10.2307/2586759}{doi:10.2307/2586759}.

\bibitem{Higman} G. Higman, \emph{Ordering by Divisibility in Abstract Algebras}, Proc. London Math. Soc. (3) \textbf{2} (1952), 326--336. \href{https://doi.org/10.1112/plms/s3-2.1.326}{doi:10.1112/plms/s3-2.1.326}.

\bibitem{Jech} T. Jech, \emph{The Axiom of Choice}, North-Holland, 1973. Example 2.3.2 and Section 7.2. \url{https://gwern.net/doc/math/1973-jech-theaxiomofchoice.pdf}.

\bibitem{Rado} R. Rado, \emph{Partial well-ordering of sets of vectors}, Mathematika \textbf{1} (1954), 89--95. \href{https://doi.org/10.1112/S0025579300000565}{doi:10.1112/S0025579300000565}.

\bibitem{Glazer} E. Glazer, answer to D. van der Zypen, \emph{Comparability of power sets and (AC)}, MathOverflow, 30 September 2025. \url{https://mathoverflow.net/a/501058}. Public mathematical communication, not a refereed article.

\bibitem{BSWY} M. Beriashvili, R. Schindler, L. Wu and L. Yu, \emph{Hamel bases and well-ordering the continuum}, Proc. Amer. Math. Soc. 146 (2018), 3565--3573. \url{https://doi.org/10.1090/proc/14010}.

\bibitem{ICFquestion} A. Karagila, \emph{Does $2^X=2^Y\Rightarrow X=Y$ imply the Axiom of Choice?}, MathOverflow, 17 November 2011, with subsequent answers and discussion. \url{https://mathoverflow.net/questions/81204}. Historical question and terminology; current open status is not inferred from the date of the page.

\bibitem{Hamkins} J. D. Hamkins, answer to A. Karagila, \emph{Does $2^X=2^Y\Rightarrow X=Y$ imply the Axiom of Choice?}, MathOverflow, 24 August 2012. \url{https://mathoverflow.net/a/105424}.

\bibitem{LT} A. Lindenbaum and A. Tarski, \emph{Communication sur les recherches de la th\'eorie des ensembles}, Comptes Rendus des S\'eances de la Soci\'et\'e des Sciences et des Lettres de Varsovie, Classe III, 19 (1926), 299--330, No.~82(A5), p.~312. Original scan: \url{https://rcin.org.pl/Content/50065}.

\bibitem{Grochow} J. A. Grochow, \emph{Polynomial-Time Axioms of Choice and Polynomial-Time Cardinality}, Theory of Computing Systems 67 (2023), 627--669. \url{https://doi.org/10.1007/s00224-023-10118-y}. Appendix A.3 reproduces the related Rubin--Rubin strict-reflection formulation.

\bibitem{KS} A. Karagila and J. Schilhan, \emph{Intermediate models and Kinna--Wagner Principles}, Proc. Amer. Math. Soc. 154 (2026), 393--403. \url{https://doi.org/10.1090/proc/17425}.

\bibitem{Blass1977} A. Blass, \emph{Ramsey's theorem in the hierarchy of choice principles}, Journal of Symbolic Logic \textbf{42} (1977), no.~3, 387--390. \href{https://doi.org/10.2307/2272866}{doi:10.2307/2272866}.

\bibitem{ErdosRado} P. Erd\H{o}s and R. Rado, \emph{A partition calculus in set theory}, Bulletin of the American Mathematical Society \textbf{62} (1956), 427--489. Theorem~39, pp.~467--469, especially clause~(iii). Original paper: \url{https://www.renyi.hu/~p_erdos/1956-02.pdf}.

\bibitem{BHT} J. E. Baumgartner, A. Hajnal and S. Todor\v{c}evi\'c, \emph{Extensions of the Erd\H{o}s--Rado theorem}, manuscript, arXiv:math/9311207 (1993). Theorem~2.1, manuscript p.~4. \url{https://arxiv.org/abs/math/9311207}.
\end{thebibliography}
